\documentclass[10pt,a4paper]{amsart}
\usepackage{amsmath,amssymb,amsthm,graphicx}
\usepackage{verbatim}
\newenvironment{reference}{\comment}{\endcomment}
%\newenvironment{reference}{}{}
\newenvironment{slogan}{\comment}{\endcomment}
\newenvironment{history}{\comment}{\endcomment}

% For commutative diagrams we use Xy-pic
\usepackage[all]{xy}

% We use 2cell for 2-commutative diagrams.
\xyoption{2cell}
% Theorem environments.
%
\theoremstyle{plain}
\newtheorem{theorem}[subsection]{Theorem}
\newtheorem{proposition}[subsection]{Proposition}
\newtheorem{lemma}[subsection]{Lemma}

\theoremstyle{definition}
\newtheorem{definition}[subsection]{Definition}
\newtheorem{example}[subsection]{Example}
\newtheorem{exercise}[subsection]{Exercise}
\newtheorem{situation}[subsection]{Situation}

\theoremstyle{remark}
\newtheorem{remark}[subsection]{Remark}
\newtheorem{remarks}[subsection]{Remarks}

\numberwithin{equation}{subsection}

% Macros
%
\def\lim{\mathop{\mathrm{lim}}\nolimits}
\def\colim{\mathop{\mathrm{colim}}\nolimits}
\def\Spec{\mathop{\mathrm{Spec}}}
\def\Hom{\mathop{\mathrm{Hom}}\nolimits}
\def\Ext{\mathop{\mathrm{Ext}}\nolimits}
\def\SheafHom{\mathop{\mathcal{H}\!\mathit{om}}\nolimits}
\def\SheafExt{\mathop{\mathcal{E}\!\mathit{xt}}\nolimits}
\def\Sch{\mathit{Sch}}
\def\Mor{\mathop{\mathrm{Mor}}\nolimits}
\def\Ob{\mathop{\mathrm{Ob}}\nolimits}
\def\Sh{\mathop{\mathit{Sh}}\nolimits}
\def\NL{\mathop{N\!L}\nolimits}
\def\CH{\mathop{\mathrm{CH}}\nolimits}
\def\proetale{{pro\text{-}\acute{e}tale}}
\def\etale{{\acute{e}tale}}
\def\QCoh{\mathit{QCoh}}
\def\Ker{\mathop{\mathrm{Ker}}}
\def\Im{\mathop{\mathrm{Im}}}
\def\Coker{\mathop{\mathrm{Coker}}}
\def\Coim{\mathop{\mathrm{Coim}}}

% Boxtimes
%
\DeclareMathSymbol{\boxtimes}{\mathbin}{AMSa}{"02}

%
% Macros for moduli stacks/spaces
%
\def\QCohstack{\mathcal{QC}\!\mathit{oh}}
\def\Cohstack{\mathcal{C}\!\mathit{oh}}
\def\Spacesstack{\mathcal{S}\!\mathit{paces}}
\def\Quotfunctor{\mathrm{Quot}}
\def\Hilbfunctor{\mathrm{Hilb}}
\def\Curvesstack{\mathcal{C}\!\mathit{urves}}
\def\Polarizedstack{\mathcal{P}\!\mathit{olarized}}
\def\Complexesstack{\mathcal{C}\!\mathit{omplexes}}
% \Pic is the operator that assigns to X its picard group, usage \Pic(X)
% \Picardstack_{X/B} denotes the Picard stack of X over B
% \Picardfunctor_{X/B} denotes the Picard functor of X over B
\def\Pic{\mathop{\mathrm{Pic}}\nolimits}
\def\Picardstack{\mathcal{P}\!\mathit{ic}}
\def\Picardfunctor{\mathrm{Pic}}
\def\Deformationcategory{\mathcal{D}\!\mathit{ef}}
\usepackage{cleveref}
\crefname{originalreference}{original reference}{original references}
\Crefname{originalreference}{Original reference}{Original references}
\makeatletter
\newlabel{remark-relative-dualizing-complex}{{0B6S}{source}{}{sourcepointer-1}{}}
\newlabel{remark-relative-dualizing-complex@cref}{{[originalreference][0][]0B6S}{[1][0][]source}{}{}{}}
\makeatother
\crefname{originalequation}{original equation}{original equations}
\Crefname{originalequation}{Original equation}{Original equations}
\makeatletter
\newlabel{equation-pre-rigid}{{0E2Q}{source}{}{sourcepointer-2}{}}
\newlabel{equation-pre-rigid@cref}{{[originalequation][0][]0E2Q}{[1][0][]source}{}{}{}}
\makeatother
\crefname{originaldefinition}{original definition}{original definitions}
\Crefname{originaldefinition}{Original definition}{Original definitions}
\makeatletter
\newlabel{definition-relative-dualizing-complex}{{0E2T}{source}{}{sourcepointer-3}{}}
\newlabel{definition-relative-dualizing-complex@cref}{{[originaldefinition][0][]0E2T}{[1][0][]source}{}{}{}}
\makeatother
\crefname{originallemma}{original lemma}{original lemmas}
\Crefname{originallemma}{Original lemma}{Original lemmas}
\makeatletter
\newlabel{lemma-properties-relative-dualizing}{{0E4L}{source}{}{sourcepointer-4}{}}
\newlabel{lemma-properties-relative-dualizing@cref}{{[originallemma][0][]0E4L}{[1][0][]source}{}{}{}}
\makeatother
\crefname{originallemma}{original lemma}{original lemmas}
\Crefname{originallemma}{Original lemma}{Original lemmas}
\makeatletter
\newlabel{lemma-twisted-inverse-image}{{0A9E}{source}{}{sourcepointer-5}{}}
\newlabel{lemma-twisted-inverse-image@cref}{{[originallemma][0][]0A9E}{[1][0][]source}{}{}{}}
\makeatother
\crefname{originallemma}{original lemma}{original lemmas}
\Crefname{originallemma}{Original lemma}{Original lemmas}
\makeatletter
\newlabel{lemma-twisted-inverse-image-closed}{{0A9X}{source}{}{sourcepointer-6}{}}
\newlabel{lemma-twisted-inverse-image-closed@cref}{{[originallemma][0][]0A9X}{[1][0][]source}{}{}{}}
\makeatother
\crefname{originallemma}{original lemma}{original lemmas}
\Crefname{originallemma}{Original lemma}{Original lemmas}
\makeatletter
\newlabel{lemma-sheaf-with-exact-support-ext}{{0A77}{source}{}{sourcepointer-7}{}}
\newlabel{lemma-sheaf-with-exact-support-ext@cref}{{[originallemma][0][]0A77}{[1][0][]source}{}{}{}}
\makeatother
\title{The relative dualizing complex of a proper flat morphism is a relative dualizing complex}
\author{The Stacks Project Authors}
\date{stacks-project:duality.tex:local-snapshot-20261012; GNU FDL 1.2 or later}
\begin{document}
\maketitle
\noindent\textit{Editorial note.} All delivered mathematical text is the author's own, copied from the pinned Stacks Project chapter file duality.tex; no statement or proof step was rewritten, reordered or removed.\par\medskip
\noindent\textit{Editorial note.} Preamble provenance: the pinned chapter file opens with an input of the shared preamble and carries no preamble text of its own, so the project's single Stacks-wide preamble staged at sources/stacks-duality/preamble.tex is excerpted verbatim under the recipe fixed by LANE\_BRIEF 22.1, as three spans, lines 8-13 (verbatim package and the reference/slogan/history environments), lines 14-18 (Xy-pic setup) and lines 154-217 (the native theorem declarations and mathematical macros). Nothing was added, removed or redefined; the wrapper class and the amsmath/amssymb/amsthm/graphicx base packages are supplied by the assembler.\par\medskip
\noindent\textit{Editorial note.} The counted statement and proof contain no cross reference, no citation and no included file, so the delivered body is self contained.\par\medskip
\noindent\textit{Editorial note.} Canonical location of the counted proposition: Stacks tag 0E2Z (sources/stacks-duality/pinned-original-source.tex line 7530).\par\medskip
\noindent\textit{Allowed standard prerequisites (editorial).}\begin{itemize}
\item The relative dualizing complex is S-perfect. Same chapter, Lemma labelled lemma-properties-relative-dualizing, cited at proof lines 7540-7542. Gives the S-perfectness required of a relative dualizing complex.
\item The right adjoint of pushforward for the diagonal. Same chapter, Lemma labelled lemma-twisted-inverse-image, cited at proof lines 7542-7546. Supplies the functor c applied to the pulled back dualizing complex.
\item The right adjoint of a closed immersion and the exact support sheaf. Same chapter, Lemmas labelled lemma-twisted-inverse-image-closed and lemma-sheaf-with-exact-support-ext, cited at proof lines 7552-7556. Identify Delta\_*(c(...)) with RSheafHom against Delta\_*O\_X.
\end{itemize}\medskip
\ifdefined\hypertarget\hypertarget{sourcepointer-1}{}\fi
\noindent\textit{Source reference.} Original reference 0B6S: sources/stacks-duality/pinned-original-source.tex, source lines 2945--2945. Same-chapter pointer into the author environment remark at source line 2945; invoked at proof line None of the counted proof. It is retained only as a navigation pointer; every load-bearing step is inside the delivered span.\par\medskip
\ifdefined\hypertarget\hypertarget{sourcepointer-2}{}\fi
\noindent\textit{Source reference.} Original equation 0E2Q: sources/stacks-duality/pinned-original-source.tex, source lines 3166--3166. Same-chapter pointer into the author environment equation at source line 3166; invoked at proof line 7546 of the counted proof. It is retained only as a navigation pointer; every load-bearing step is inside the delivered span.\par\medskip
\ifdefined\hypertarget\hypertarget{sourcepointer-3}{}\fi
\noindent\textit{Source reference.} Original definition 0E2T: sources/stacks-duality/pinned-original-source.tex, source lines 7234--7234. Same-chapter pointer into the author environment definition at source line 7234; invoked at proof line None of the counted proof. It is retained only as a navigation pointer; every load-bearing step is inside the delivered span.\par\medskip
\ifdefined\hypertarget\hypertarget{sourcepointer-4}{}\fi
\noindent\textit{Source reference.} Original lemma 0E4L: sources/stacks-duality/pinned-original-source.tex, source lines 3028--3028. Same-chapter pointer into the author environment lemma at source line 3028; invoked at proof line 7541 of the counted proof. It is retained only as a navigation pointer; every load-bearing step is inside the delivered span.\par\medskip
\ifdefined\hypertarget\hypertarget{sourcepointer-5}{}\fi
\noindent\textit{Source reference.} Original lemma 0A9E: sources/stacks-duality/pinned-original-source.tex, source lines 653--653. Same-chapter pointer into the author environment lemma at source line 653; invoked at proof line 7544 of the counted proof. It is retained only as a navigation pointer; every load-bearing step is inside the delivered span.\par\medskip
\ifdefined\hypertarget\hypertarget{sourcepointer-6}{}\fi
\noindent\textit{Source reference.} Original lemma 0A9X: sources/stacks-duality/pinned-original-source.tex, source lines 2467--2467. Same-chapter pointer into the author environment lemma at source line 2467; invoked at proof line 7555 of the counted proof. It is retained only as a navigation pointer; every load-bearing step is inside the delivered span.\par\medskip
\ifdefined\hypertarget\hypertarget{sourcepointer-7}{}\fi
\noindent\textit{Source reference.} Original lemma 0A77: sources/stacks-duality/pinned-original-source.tex, source lines 2339--2339. Same-chapter pointer into the author environment lemma at source line 2339; invoked at proof line 7556 of the counted proof. It is retained only as a navigation pointer; every load-bearing step is inside the delivered span.\par\medskip
\begin{lemma}
\label{lemma-flat-proper-relative-dualizing}
Let $S$ be a quasi-compact and quasi-separated scheme.
Let $f : X \to S$ be a proper, flat morphism of finite presentation.
The relative dualizing complex $\omega_{X/S}^\bullet$ of
Remark \ref{remark-relative-dualizing-complex}
together with (\ref{equation-pre-rigid}) is a relative
dualizing complex in the sense of
Definition \ref{definition-relative-dualizing-complex}.
\end{lemma}

\begin{proof}
In Lemma \ref{lemma-properties-relative-dualizing} we proved that
$\omega_{X/S}^\bullet$ is $S$-perfect.
Let $c$ be the right adjoint of
Lemma \ref{lemma-twisted-inverse-image}
for the diagonal $\Delta : X \to X \times_S X$.
Then we can apply $\Delta_*$ to (\ref{equation-pre-rigid})
to get an isomorphism
$$
\Delta_*\mathcal{O}_X \to
\Delta_*(c(L\text{pr}_1^*\omega_{X/S}^\bullet)) =
R\SheafHom_{\mathcal{O}_{X \times_S X}}(
\Delta_*\mathcal{O}_X, L\text{pr}_1^*\omega_{X/S}^\bullet)
$$
The equality holds by
Lemmas \ref{lemma-twisted-inverse-image-closed} and
\ref{lemma-sheaf-with-exact-support-ext}.
This finishes the proof.
\end{proof}

\end{document}
